\section{\ENG{Observability} \& \ENG{Governance}}
\label{sec:observability_governance}

Το \ENG{monitoring} της πλατφόρμας στηρίζεται πρωτίστως σε \ENG{structured logging}. Στην πλευρά του \ENG{engine}, το \texttt{\ENG{logback-spring.xml}} διαμορφώνει \ENG{JSON logs}, ενώ οι \texttt{\ENG{CorrelationIdFilter}} και \texttt{\ENG{RequestLogFilter}} προσθέτουν αναγνωριστικό συσχέτισης και συνοπτική πληροφορία για τις \ENG{HTTP} αιτήσεις. Η ύπαρξη του \ENG{X-Correlation-Id} στο φίλτρο καθιστά εφικτή την παρακολούθηση μιας αλυσίδας αιτήσεων μέσα από διαφορετικά επίπεδα της πλατφόρμας, ακόμη κι αν δεν υπάρχει πλήρες \ENG{distributed tracing} σύστημα.

Στην πλευρά του \ENG{player}, η χρήση \ENG{pino} και \ENG{pino-http} παρέχει συνεπή καταγραφή για τις αιτήσεις και για τις κλήσεις προς το \ENG{engine}, ιδίως μέσω του \texttt{\ENG{ObservedEngineClient}}. Παράλληλα, στο επίπεδο της κατάστασης εκτέλεσης υπάρχει περιορισμένη αλλά σαφής μετρική παρατηρησιμότητα: ο \texttt{\ENG{ObservedRuntimeStateStore}} τυλίγει το κανονικό αποθετήριο κατάστασης και εκθέτει γεγονότα λειτουργίας προς \ENG{Micrometer}. Το σημείο αυτό είναι αρκετό για να στηρίξει ακαδημαϊκά τον ισχυρισμό ότι η πλατφόρμα διαθέτει βασικό \ENG{instrumentation} γύρω από την κρίσιμη κατάσταση εκτέλεσης, αλλά όχι για να αποδοθεί πλήρης υποδομή μετρικών με πλούσια \ENG{dashboards}.

Η λειτουργική διαχειρίση ολοκληρώνεται από τα \ENG{CI workflows} των αποθετηρίων και από τα \ENG{scripts} ελέγχου του κεντρικού περιβάλλοντος. Οι διαφορετικοί τρόποι εκτέλεσης δοκιμών, η ύπαρξη \ENG{mock compose} σεναρίων για ελέγχους \ENG{pull request} και η συστηματική επαλήθευση \ENG{readiness} πριν από τα \ENG{smoke checks} δείχνουν ότι η πλατφόρμα ελέγεχται σε όλα τα στάδια ανάπτυξης και λειτουργίας, από τον κώδικα μέχρι την παραγωγή.

Το Σχήμα~\ref{fig:observability_governance} συνοψίζει τα κανάλια \ENG{logs}, συσχέτισης και βασικών μετρικών, ενώ ο Πίνακας~\ref{tab:operational_controls_and_checks} αποδίδει συνολικά τα κύρια \ENG{operational controls}. Η ανάλυση αυτή αρκεί για να τεκμηριώσει υλοποιημένη παρατηρησιμότητα και λειτουργική πειθαρχία, χωρίς να υπερβαίνει τα όρια αυτού που πράγματι υπάρχει στον κώδικα και στην υποδομή.

\begin{figure}[H]
  \centering
\begin{chapterfigurebox}
  \begin{tikzpicture}[
      font=\footnotesize,
      node distance=1.6cm and 2.2cm,
      box/.style={draw, rounded corners, fill=gray!5, align=center, text width=3.1cm, minimum height=1.05cm},
      emph/.style={draw, rounded corners, fill=gray!12, align=center, text width=3.4cm, minimum height=1.05cm, font=\footnotesize\bfseries},
      arrow/.style={-Latex, thick}
    ]
    \node[box] (enginelogs) {\ENG{Engine JSON logs}};
    \node[box, below=0.8cm of enginelogs] (playerlogs) {\ENG{Player pino logs}};
    \node[box, below=0.8cm of playerlogs] (containerlogs) {\ENG{Container logs}\\\ENG{via gelf}};

    \node[emph, right=2.3cm of playerlogs] (logstash) {\ENG{Logstash}};
    \node[emph, right=2.3cm of logstash] (elastic) {\ENG{Elasticsearch}};
    \node[emph, right=2.3cm of elastic] (kibana) {\ENG{Kibana}};

    \node[box, above=0.9cm of logstash] (corr) {\texttt{\ENG{CorrelationIdFilter}}};
    \node[box, below=1.0cm of elastic] (metrics) {\texttt{\ENG{ObservedRuntimeStateStore}}\\+\ENG{Micrometer}};

    \draw[arrow] (enginelogs) -- (logstash);
    \draw[arrow] (playerlogs) -- (logstash);
    \draw[arrow] (containerlogs) -- (logstash);
    \draw[arrow] (logstash) -- (elastic);
    \draw[arrow] (elastic) -- (kibana);
    \draw[arrow] (corr) -- (enginelogs);
    \draw[arrow] (metrics) -- (elastic);
  \end{tikzpicture}
\end{chapterfigurebox}
  \caption{Υλοποιημένα κανάλια παρατηρησιμότητας της πλατφόρμας: \ENG{structured logs}, συσχέτιση αιτήσεων και βασική μετρική κάλυψη της κατάστασης εκτέλεσης.}
  \label{fig:observability_governance}
\end{figure}


\begin{table}[H]
\centering
\caption{Υλοποιημένοι λειτουργικοί έλεγχοι, σήματα παρατηρησιμότητας και σημεία επιχειρησιακής επαλήθευσης}
\label{tab:operational_controls_and_checks}
\begingroup
\footnotesize
\setlength{\tabcolsep}{3pt}
\renewcommand{\arraystretch}{1.20}
\begin{adjustbox}{center,max width=\textwidth}
\begin{tabular}{|p{3.4cm}|p{4.2cm}|p{5.2cm}|}
\hline
\ENG{Μηχανισμός} & Υλοποιημένο \ENG{artifact} & Λειτουργικός ρόλος \\
\hline
\ENG{Structured engine logging} & \texttt{\ENG{logback-spring.xml}}, \texttt{\ENG{RequestLogFilter}} & Συνεπής \ENG{JSON} καταγραφή αιτήσεων και γεγονότων πυρήνα \\
\hline
\ENG{Correlation IDs} & \texttt{\ENG{CorrelationIdFilter}} & Συσχέτιση αιτήσεων σε πολλαπλά επίπεδα του συστήματος \\
\hline
\ENG{Player logging} & \ENG{pino}, \ENG{pino-http}, \texttt{\ENG{ObservedEngineClient}} & Καταγραφή αιτήσεων του \ENG{player} και κλήσεων προς το \ENG{engine} \\
\hline
Μετρικές κατάστασης εκτέλεσης & \texttt{\ENG{ObservedRuntimeStateStore}} με \ENG{Micrometer} & Βασική κάλυψη μετρικής παρακολούθησης για λειτουργίες του αποθετηρίου κατάστασης \\
\hline
Κεντρική συγκέντρωση \ENG{logs} & \ENG{gelf driver}, \texttt{\ENG{logstash.conf}}, \ENG{Elasticsearch}, \ENG{Kibana} & Αναζήτηση και επιθεώρηση \ENG{logs} σε επίπεδο πλατφόρμας \\
\hline
Έλεγχος παραμέτρων περιβάλλοντος & \texttt{\ENG{validate-env.sh}} & Έγκαιρη ανίχνευση ελλιπών ή λανθασμένων μεταβλητών \ENG{.env} \\
\hline
\ENG{Smoke checks} σύνθεσης & \texttt{\ENG{smoke-check.sh}} & Επαλήθευση \ENG{readiness} και προσβασιμότητας κρίσιμων υπηρεσιών \\
\hline
\ENG{CI workflows} & \ENG{GitHub Actions} στα αποθετήρια και \ENG{mock compose mode} στην \ENG{infra} σύνθεση & Επαναλήψιμη επαλήθευση \ENG{build}, \ENG{tests} και βασικής λειτουργικότητας \\
\hline
\end{tabular}
\end{adjustbox}
\endgroup

\end{table}
